\chapter{Arquitectura Editorial}

\begin{planningbox}{Objetivo del suplemento}
El futuro cuaderno de problemas contextualizados se concibe como un complemento de alta
dificultad del libro principal. Su funcion no es repetir la practica basica, sino obligar a
reconocer tecnicas, traducir datos de contexto a lenguaje matematico, justificar elecciones y
cerrar cada tarea con interpretacion.
\end{planningbox}

\begin{databox}{Magnitud del proyecto}
Con la estructura actual del libro, el suplemento queda proyectado sobre
\textbf{\SupplementChapterCount\ capitulos troncales} y \textbf{\SupplementSectionCount\
secciones activas}. La primera version operativa de este proyecto editorial recomienda:
\begin{itemize}
  \item un problema largo por seccion;
  \item un reto integrador por capitulo;
  \item un cierre final con tres tareas transversales;
  \item solucion completa en la misma maqueta de trabajo para facilitar la redaccion docente.
\end{itemize}
\end{databox}

\section{Criterios de diseno}

\begin{itemize}
  \item Cada problema debe tener un contexto realista, pero la matematica no puede quedar
  escondida: el alumnado debe reconocer la estructura del apartado de origen.
  \item Debe haber un recurso visual en cada pieza: tabla, recta numerica, plano, croquis,
  semiplanos, grafica, diagrama vectorial o esquema de datos.
  \item La dificultad se graduara en tres capas: lectura del contexto, modelizacion y cierre
  matematico.
  \item Las soluciones deben explicar por que se elige un metodo y no solo mostrar la cuenta.
\end{itemize}

\section{Mapa por secciones}

\setlength{\LTpre}{0pt}
\setlength{\LTpost}{0pt}
\small
\begin{longtable}{@{}>{\raggedright\arraybackslash}p{0.10\textwidth}>{\raggedright\arraybackslash}p{0.21\textwidth}>{\raggedright\arraybackslash}p{0.28\textwidth}>{\raggedright\arraybackslash}p{0.19\textwidth}>{\raggedright\arraybackslash}p{0.08\textwidth}@{}}
\toprule
\textbf{ID} & \textbf{Seccion} & \textbf{Problema contextualizado propuesto} & \textbf{Recurso visual} & \textbf{Nivel}\\
\midrule
\endfirsthead
\toprule
\textbf{ID} & \textbf{Seccion} & \textbf{Problema contextualizado propuesto} & \textbf{Recurso visual} & \textbf{Nivel}\\
\midrule
\endhead
\multicolumn{5}{l}{\textbf{C01 --- Numeros reales y herramientas numericas}}\\
C01.S01 & Clasificacion de numeros y conjuntos numericos & Auditoria de registros numericos de un laboratorio escolar & Tabla de magnitudes y unidades & Alta\\
C01.S02 & Decimales y fraccion generatriz & Consumos periodicos y reparto exacto de costes & Tabla de magnitudes y unidades & Alta\\
C01.S03 & Intervalos, union e interseccion & Ventanas de seguridad y franjas de operacion & Recta numerica o banda de tolerancias & Alta\\
C01.S04 & Valor absoluto como distancia & Control de tolerancias respecto a un valor objetivo & Recta numerica o banda de tolerancias & Alta\\
C01.S05 & Potencias y radicales & Escalado de superficies, volumenes o intensidades & Tabla de magnitudes y unidades & Alta\\
C01.S06 & Racionalizacion & Simplificacion de formulas instrumentales con radicales & Tabla de magnitudes y unidades & Alta\\
C01.S07 & Logaritmos y propiedades operativas & Lectura de escalas de sonido, pH o intensidad & Tabla de magnitudes y unidades & Alta\\
\multicolumn{5}{l}{\textbf{C02 --- Polinomios y fracciones algebraicas}}\\
C02.S01 & Valor numerico y lectura de polinomios & Modelo de costes dependiente de un parametro operativo & Tabla de datos parametrica & Alta\\
C02.S02 & Productos notables y operaciones con polinomios & Ampliacion modular de superficies y margenes de error & Tabla de datos parametrica & Alta\\
C02.S03 & Division de polinomios & Reparto por lotes con interpretacion de cociente y resto & Tabla de evaluacion y caja algebraica & Alta\\
C02.S04 & Teorema del resto, del factor y parametros & Calibracion de parametros mediante condiciones de divisibilidad & Tabla de evaluacion y caja algebraica & Muy alta\\
C02.S05 & Raices y factorizacion & Busqueda de umbrales de equilibrio en un modelo algebraico & Tabla de evaluacion y caja algebraica & Alta\\
C02.S06 & Fracciones algebraicas: dominio y simplificacion & Indice tecnico con restricciones de funcionamiento & Tabla de datos parametrica & Alta\\
C02.S07 & Suma y resta de fracciones algebraicas & Combinacion de tasas con distintos denominadores operativos & Tabla de datos parametrica & Alta\\
C02.S08 & Producto, cociente y operaciones combinadas & Encadenado de rendimientos y tiempos relativos & Tabla de datos parametrica & Alta\\
C02.S09 & Ecuaciones con fracciones algebraicas & Ajuste de tiempos medios con restricciones fisicas & Tabla de datos parametrica & Alta\\
\multicolumn{5}{l}{\textbf{C03 --- Ecuaciones y sistemas}}\\
C03.S01 & Ecuaciones racionales & Comparacion de rendimientos con tiempos de espera & Esquema de restricciones y pasos & Alta\\
C03.S02 & Ecuaciones irracionales & Diseno de rampas o diagonales con medida exigida & Esquema de restricciones y pasos & Alta\\
C03.S03 & Ecuaciones polinomicas reducibles & Modelo de produccion con variables encadenadas & Esquema de restricciones y pasos & Alta\\
C03.S04 & Ecuaciones exponenciales & Crecimiento de usuarios o capital con metas cerradas & Esquema de restricciones y pasos & Alta\\
C03.S05 & Ecuaciones logaritmicas & Recuperacion de una variable a partir de una escala logaritmica & Esquema de restricciones y pasos & Alta\\
C03.S06 & Eleccion de metodo en ecuaciones mixtas & Diagnostico de estrategia en una incidencia tecnica & Esquema de restricciones y pasos & Alta\\
C03.S07 & Sistemas lineales 2x2 & Venta de entradas o mezclas con dos restricciones & Tabla de datos o esquema de relaciones & Alta\\
C03.S08 & Sistemas lineales 3x3 y clasificacion & Auditoria de registros numericos de un laboratorio escolar & Tabla de datos o esquema de relaciones & Alta\\
C03.S09 & Problemas de planteamiento con sistemas & Problema de precios, edades o mezclas con traduccion completa & Tabla de datos o esquema de relaciones & Alta\\
C03.S10 & Sistemas no lineales & Interseccion de restricciones geometricas o economicas & Tabla de datos o esquema de relaciones & Alta\\
C03.S11 & Sistemas exponenciales y logaritmicos & Crecimiento de usuarios o capital con metas cerradas & Tabla de datos o esquema de relaciones & Muy alta\\
\multicolumn{5}{l}{\textbf{C04 --- Inecuaciones}}\\
C04.S01 & Inecuaciones lineales & Condicion minima de rentabilidad o seguridad & Tabla de signos y recta numerica & Alta\\
C04.S02 & Inecuaciones de segundo grado & Intervalo viable en un modelo cuadratico & Tabla de signos y recta numerica & Alta\\
C04.S03 & Inecuaciones polinomicas & Estudio de signos de un indice de estabilidad & Tabla de signos y recta numerica & Alta\\
C04.S04 & Inecuaciones racionales & Comparacion de rendimientos con tiempos de espera & Tabla de signos y recta numerica & Alta\\
C04.S05 & Sistemas de inecuaciones en una variable & Compatibilidad simultanea de varias normas operativas & Tabla de signos y recta numerica & Alta\\
C04.S06 & Inecuaciones de dos variables y semiplanos & Semiplanos de produccion con restricciones de recursos & Semiplanos y region factible & Alta\\
C04.S07 & Sistemas de inecuaciones en el plano & Region factible para un plan de fabricacion & Semiplanos y region factible & Alta\\
\multicolumn{5}{l}{\textbf{C05 --- Trigonometria}}\\
C05.S01 & Grados y radianes & Giros de camaras y brujulas en doble sistema de medida & Dibujo vectorial o plano cartesiano & Alta\\
C05.S02 & Razones trigonometricas en triangulos y poligonos & Medicion indirecta de alturas y diagonales & Croquis trigonometrico y tabla de datos & Alta\\
C05.S03 & Angulos notables y reduccion al primer cuadrante & Orientacion de paneles con cuadrantes y angulos de referencia & Dibujo vectorial o plano cartesiano & Alta\\
C05.S04 & Calculo de razones a partir de una razon conocida & Reconstruccion de pendientes y longitudes desde una sola razon & Dibujo vectorial o plano cartesiano & Alta\\
C05.S05 & Identidades y transformaciones trigonometricas & Simplificacion de señales periodicas equivalentes & Croquis trigonometrico y tabla de datos & Alta\\
C05.S06 & Ecuaciones trigonometricas & Sincronizacion de ciclos y ventanas de fase & Croquis trigonometrico y tabla de datos & Alta\\
C05.S07 & Resolucion de triangulos & Topografia de parcelas y trazados & Croquis trigonometrico y tabla de datos & Alta\\
C05.S08 & Problemas contextualizados de trigonometria & Alturas, sombras y navegacion con datos reales & Croquis trigonometrico y tabla de datos & Alta\\
\multicolumn{5}{l}{\textbf{C06 --- Vectores}}\\
C06.S01 & Bases y dependencia en R2 & Composicion de desplazamientos sobre dos direcciones base & Dibujo vectorial o plano cartesiano & Alta\\
C06.S02 & Operaciones con vectores & Ruta compuesta de un dron o robot movil & Dibujo vectorial o plano cartesiano & Alta\\
C06.S03 & Modulo, argumento y vectores unitarios & Vector de viento dominante y correccion angular & Dibujo vectorial o plano cartesiano & Alta\\
C06.S04 & Producto escalar y angulo entre vectores & Trabajo mecanico o compatibilidad de trayectorias & Dibujo vectorial o plano cartesiano & Alta\\
C06.S05 & Paralelismo, perpendicularidad y alineacion & Alineacion de sensores y piezas & Dibujo vectorial o plano cartesiano & Alta\\
C06.S06 & Construcciones con vectores y paralelogramos & Diseno de una estructura con diagonales y lados opuestos & Dibujo vectorial o plano cartesiano & Alta\\
\multicolumn{5}{l}{\textbf{C07 --- Geometria analitica y lugares geometricos}}\\
C07.S01 & Ecuaciones de la recta y vectores asociados & Trazado de una calle o conduccion sobre plano & Plano cartesiano con TikZ & Alta+\\
C07.S02 & Problemas con parametros en rectas & Calibracion geometrica para cumplir una condicion & Plano cartesiano con TikZ & Muy alta\\
C07.S03 & Simetrias respecto de punto y recta & Reflejo de trayectorias respecto a ejes urbanos & Plano cartesiano con TikZ & Alta+\\
C07.S04 & Posicion relativa de rectas & Cruce, paralelismo o coincidencia de recorridos & Plano cartesiano con TikZ & Alta+\\
C07.S05 & Rectas especiales por un punto & Recta de servicio perpendicular o paralela desde un punto & Plano cartesiano con TikZ & Alta+\\
C07.S06 & Distancias y angulos entre rectas & Separacion minima entre infraestructuras & Plano cartesiano con TikZ & Alta+\\
C07.S07 & Poligonos y areas con coordenadas & Area de un recinto medido por coordenadas & Plano cartesiano con TikZ & Alta+\\
C07.S08 & Medianas, alturas y mediatrices & Centros y rectas notables de una parcela triangular & Dibujo vectorial o plano cartesiano & Alta+\\
C07.S09 & Bisectrices, ortocentro y lugares geometricos & Punto de control equidistante y geometria de triangulos & Plano cartesiano con TikZ & Alta+\\
C07.S10 & Condiciones geometricas combinadas & Diseno de una pieza con varias restricciones simultaneas & Dibujo vectorial o plano cartesiano & Muy alta\\
\multicolumn{5}{l}{\textbf{C08 --- Limites, continuidad y asintotas}}\\
C08.S01 & Dominio de funciones & Modelo fisico con restricciones de existencia & Tabla de valores y grafica de apoyo & Alta+\\
C08.S02 & Lectura de graficas & Analisis visual de un proceso registrado & Grafica con PGFPlots & Alta+\\
C08.S03 & Limites basicos y laterales & Comportamiento cercano a un umbral de medida & Tabla de valores y grafica de apoyo & Alta+\\
C08.S04 & Limites en el infinito & Estudio de evolucion a largo plazo & Tabla de valores y grafica de apoyo & Alta+\\
C08.S05 & Indeterminaciones por factorizacion y racionalizacion & Simplificacion de formulas instrumentales con radicales & Tabla de valores y grafica de apoyo & Alta+\\
C08.S06 & Asintotas & Modelo con barreras y tendencia lineal & Grafica con PGFPlots & Alta+\\
C08.S07 & Continuidad y discontinuidades & Transicion de un sistema entre distintos regimenes & Tabla de valores y grafica de apoyo & Alta+\\
C08.S08 & Funciones a trozos y representacion & Tarifa escalonada o politica de precios por tramos & Grafica con PGFPlots & Alta+\\
C08.S09 & Continuidad con parametros & Calibracion de una funcion para evitar saltos & Tabla de valores y grafica de apoyo & Muy alta\\
C08.S10 & Modelos de evolucion a largo plazo & Modelo de saturacion o enfriamiento & Tabla de valores y grafica de apoyo & Alta+\\
C08.S11 & Ampliacion y estilo EBAU en limites y continuidad & Item de analisis global con lectura estructural & Tabla de valores y grafica de apoyo & Muy alta\\
\multicolumn{5}{l}{\textbf{C09 --- Derivadas y aplicaciones}}\\
C09.S01 & Reglas de derivacion & Ritmo instantaneo de una magnitud compuesta & Tabla de valores y grafica de apoyo & Alta+\\
C09.S02 & Monotonia, extremos, curvatura e inflexion & Estudio cualitativo de una funcion de rendimiento & Tabla de valores y grafica de apoyo & Alta+\\
C09.S03 & Recta tangente y recta normal en un punto & Aproximacion lineal local de una trayectoria & Tabla de valores y grafica de apoyo & Alta+\\
C09.S04 & Tangente con pendiente dada o paralela a una recta & Busqueda de tangencias con condicion de direccion & Tabla de valores y grafica de apoyo & Alta+\\
C09.S05 & Problemas con parametros en derivadas & Ajuste de parametros para imponer extremos o tangencias & Tabla de valores y grafica de apoyo & Muy alta\\
C09.S06 & Derivabilidad de funciones a trozos & Tarifa escalonada o politica de precios por tramos & Tabla de valores y grafica de apoyo & Alta+\\
C09.S07 & Optimizacion & Diseno optimo de coste, area o beneficio & Tabla de valores y grafica de apoyo & Muy alta\\
C09.S08 & Ampliacion y selectividad en derivadas & Estudio global tipo examen externo & Tabla de valores y grafica de apoyo & Muy alta\\
\multicolumn{5}{l}{\textbf{C10 --- Estadistica, probabilidad e inferencia}}\\
C10.S01 & Datos bidimensionales y tablas de contingencia & Control de habitos de estudio mediante una tabla de contingencia & Tabla de contingencia & Alta\\
C10.S02 & Nubes de puntos, correlacion y causalidad & Auditoria de una relacion cientifica aparente & Nube de puntos, ajuste y grafico de residuos & Alta\\
C10.S03 & Regresion lineal y cuadratica & Calibracion y prediccion responsable de un sensor & Nube de puntos, ajuste y grafico de residuos & Alta\\
C10.S04 & Muestreo, inferencia y tecnologia & Encuesta ambiental con riesgo de sesgo & Tabla de muestra por estratos & Alta\\
C10.S05 & Experimentos aleatorios y algebra de sucesos & Fiabilidad conjunta de dos componentes & Diagrama de Venn & Alta\\
C10.S06 & Regla de Laplace y tecnicas de recuento & Seleccion de equipos de laboratorio & Arbol de recuento & Alta\\
C10.S07 & Probabilidad condicionada, total y Bayes & Sistema de alerta con falsos positivos & Diagrama de arbol de probabilidades & Alta\\
\bottomrule
\end{longtable}


\chapter{Banco Semilla de Problemas Modelo}

\begin{planningbox}{Como leer este banco}
Los diez problemas siguientes no pretenden agotar el suplemento. Funcionan como semilla
editorial: fijan tono, densidad, extension de solucion, tipo de grafico y nivel de exigencia.
Cada uno puede inspirar despues varios problemas paralelos dentro del mismo capitulo.
\end{planningbox}

\section{C01. Escala logaritmica de ruido}

\begin{contextproblem}{Biblioteca, patio y escala de decibelios}
La biblioteca de un centro mide el nivel sonoro mediante la expresion
\[
L=10\log\left(\frac{I}{I_0}\right),
\]
donde \(L\) es el nivel en decibelios, \(I_0\) es la intensidad de referencia y \(\log\)
representa el logaritmo decimal.

En una medicion se registran dos zonas:
\begin{enumerate}
  \item sala de estudio: \(45\) dB;
  \item patio cubierto: \(72\) dB.
\end{enumerate}

Se pide:
\begin{enumerate}
  \item comparar la intensidad del patio con la de la sala;
  \item fijar el mayor nivel admisible en una tercera zona para que su intensidad no sea mas de
  mil veces la de la sala de estudio.
\end{enumerate}
\end{contextproblem}

\begin{databox}{Tabla de mediciones}
\begin{center}
\begin{tabular}{lcc}
\toprule
Zona & Nivel registrado & Magnitud a estudiar \\
\midrule
Sala de estudio & \(45\) dB & Intensidad base \\
Patio cubierto & \(72\) dB & Factor de comparacion \\
Zona de lectura rapida & ? & Limite de seguridad \\
\bottomrule
\end{tabular}
\end{center}
\end{databox}

\begin{shortanswer}{C01 modelo}
El patio tiene una intensidad aproximadamente \(10^{2.7}\approx 501\) veces mayor que la sala.
Para no superar un factor \(1000\), el tercer espacio no debe pasar de \(75\) dB.
\end{shortanswer}

\begin{fullsolution}{C01 modelo}
\textbf{Paso 1. Relacionar la diferencia de decibelios con el cociente de intensidades.}
Si \(L_1=45\) y \(L_2=72\), entonces
\[
L_2-L_1=10\log\left(\frac{I_2}{I_1}\right).
\]
Por tanto,
\[
27=10\log\left(\frac{I_2}{I_1}\right)
\Longrightarrow
\log\left(\frac{I_2}{I_1}\right)=2.7
\Longrightarrow
\frac{I_2}{I_1}=10^{2.7}\approx 501.19.
\]
El patio multiplica aproximadamente por \(501\) la intensidad de la sala.

\textbf{Paso 2. Traducir el limite de intensidad a decibelios.}
Para el limite de mil veces:
\[
\frac{I}{I_1}\leq 1000=10^3.
\]
Entonces
\[
L-L_1=10\log\left(\frac{I}{I_1}\right)\leq 10\log(10^3)=30.
\]
Como \(L_1=45\),
\[
L\leq 45+30=75.
\]
El nivel maximo recomendable para la tercera zona es \(75\) dB.

\textbf{Paso 3. Comprobar.} La diferencia respecto a la sala es \(75-45=30\) dB y
\(10^{30/10}=10^3=1000\). Por tanto, el valor limite cumple exactamente la condicion.
\end{fullsolution}

\section{C02. Produccion por lotes y polinomios}

\begin{contextproblem}{Margen de un taller segun el numero de lotes}
El margen semanal \(M(n)\), en centenas de euros, de un taller que fabrica \(n\) lotes se
modeliza por
\[
M(n)=n^3-4n^2-n+4.
\]
Se pide:
\begin{enumerate}
  \item factorizar el polinomio para detectar niveles de produccion de equilibrio;
  \item dividir \(M(n)\) entre \(n-2\) e interpretar el resto;
  \item decidir si el margen se anula en \(n=2\).
\end{enumerate}
\end{contextproblem}

\begin{shortanswer}{C02 modelo}
\[
M(n)=(n-1)(n+1)(n-4),\qquad
M(n)=(n-2)(n^2-2n-5)-6.
\]
Los ceros algebraicos son \(n=-1\), \(n=1\) y \(n=4\). Si \(n\) representa un numero de lotes,
\(n=-1\) se descarta: los niveles de produccion de equilibrio son \(n=1\) y \(n=4\). En
\(n=2\), \(M(2)=-6\).
\end{shortanswer}

\begin{fullsolution}{C02 modelo}
\textbf{Paso 1. Factorizar para localizar los ceros.}
Agrupando,
\[
M(n)=n^2(n-4)-1(n-4)=(n^2-1)(n-4)=(n-1)(n+1)(n-4).
\]
Luego los ceros algebraicos del modelo son
\[
n=1,\qquad n=-1,\qquad n=4.
\]
Como el numero de lotes debe ser un entero no negativo, \(n=-1\) no tiene sentido en el
contexto. Los niveles de equilibrio fisicamente admisibles son \(n=1\) y \(n=4\).

\textbf{Paso 2. Dividir entre \(n-2\).}
Dividiendo entre \(n-2\),
\[
\begin{array}{r|rrrr}
2 & 1 & -4 & -1 & 4\\
  &   & 2 & -4 & -10\\
\hline
  & 1 & -2 & -5 & -6
\end{array}
\]
por lo que
\[
M(n)=(n-2)(n^2-2n-5)-6.
\]
\textbf{Paso 3. Interpretar y comprobar el resto.} El teorema del resto afirma que ese resto es
\(M(2)\). La sustitucion directa confirma
\[
M(2)=2^3-4\cdot2^2-2+4=8-16-2+4=-6.
\]
Por tanto, con \(2\) lotes el margen es \(-6\) centenas de euros, es decir, una perdida de
\(600\) euros; no se alcanza el equilibrio.
\end{fullsolution}

\section{C03. Suscripciones y sistema lineal 3x3}

\begin{contextproblem}{Campana de matriculas en tres planes}
Una academia ofrece tres planes: basico, estandar y premium. En una campana se venden
\(120\) suscripciones y se recauda un total de \(3030\) euros. Cada plan cuesta:
\[
18\text{ euros},\qquad 25\text{ euros},\qquad 40\text{ euros}.
\]
Ademas, el numero de planes basicos duplica al de premium.

Calcula cuantas suscripciones de cada tipo se vendieron.
\end{contextproblem}

\begin{shortanswer}{C03 modelo}
Se vendieron \(60\) planes basicos, \(30\) estandar y \(30\) premium.
\end{shortanswer}

\begin{fullsolution}{C03 modelo}
\textbf{Paso 1. Definir las incognitas y traducir los datos.}
Sean \(b,s,p\) los planes basico, estandar y premium. El sistema es
\[
\begin{cases}
b+s+p=120,\\
18b+25s+40p=3030,\\
b=2p.
\end{cases}
\]
La primera ecuacion cuenta suscripciones, la segunda representa los ingresos y la tercera
traduce que los planes basicos duplican a los premium.

\textbf{Paso 2. Reducir el sistema a una incognita.}
Sustituyendo \(b=2p\) en la primera ecuacion,
\[
2p+s+p=120 \Longrightarrow s=120-3p.
\]
En la ecuacion de ingresos:
\[
18(2p)+25(120-3p)+40p=3030.
\]
Desarrollando,
\[
36p+3000-75p+40p=3030
\Longrightarrow p+3000=3030
\Longrightarrow p=30.
\]
\textbf{Paso 3. Recuperar las demas incognitas.}
Entonces
\[
b=2p=60,\qquad s=120-3p=30.
\]
La distribucion final es:
\[
(b,s,p)=(60,30,30).
\]
\textbf{Paso 4. Comprobar en el contexto.}
\[
60+30+30=120,
\]
\[
18\cdot60+25\cdot30+40\cdot30=1080+750+1200=3030,
\qquad 60=2\cdot30.
\]
Las tres condiciones se cumplen y las cantidades son enteras no negativas.
\end{fullsolution}

\section{C04. Region factible y beneficio}

\begin{contextproblem}{Plan de fabricacion con dos productos}
Una cooperativa produce cuadernos \(x\) y archivadores \(y\). Las restricciones semanales son:
\[
\begin{cases}
x+2y\leq 140,\\
3x+2y\leq 240,\\
x\geq 20,\\
y\geq 10.
\end{cases}
\]
El beneficio es
\[
B(x,y)=18x+25y.
\]
Se pide representar la region factible, hallar sus vertices y decidir en que punto se maximiza
el beneficio.
\end{contextproblem}

\begin{databox}{Croquis de la region factible}
\begin{center}
\begin{tikzpicture}[scale=0.08]
  \draw[->] (0,0) -- (90,0) node[right] {\(x\)};
  \draw[->] (0,0) -- (0,80) node[above] {\(y\)};
  \draw[thick] (0,70) -- (140,0);
  \draw[thick] (0,120) -- (80,0);
  \draw[dashed] (20,0) -- (20,80);
  \draw[dashed] (0,10) -- (90,10);
  \fill[gray!25] (20,10) -- (20,60) -- (50,45) -- (73.333,10) -- cycle;
  \node at (58,57) {\scriptsize region viable};
  \fill (20,10) circle (1.5pt);
  \fill (20,60) circle (1.5pt);
  \fill (50,45) circle (1.5pt);
  \fill (73.333,10) circle (1.5pt);
\end{tikzpicture}
\end{center}
\end{databox}

\begin{shortanswer}{C04 modelo}
Los vertices son \((20,10)\), \((20,60)\), \((50,45)\) y
\(\left(\frac{220}{3},10\right)\). El beneficio maximo es \(2025\) euros y se alcanza en
\((50,45)\).
\end{shortanswer}

\begin{fullsolution}{C04 modelo}
\textbf{Paso 1. Dibujar las fronteras y elegir los semiplanos.}
Las fronteras son \(x+2y=140\), \(3x+2y=240\), \(x=20\) e \(y=10\). El punto
\((30,20)\) satisface estrictamente las cuatro desigualdades, de modo que permite identificar el lado valido
de cada frontera.

\textbf{Paso 2. Calcular los vertices consecutivos.}
La interseccion de \(x=20\) e \(y=10\) es \((20,10)\). Para \(x=20\) en la primera frontera,
\[
20+2y=140 \Longrightarrow y=60,
\]
lo que da \((20,60)\). La interseccion de
\[
x+2y=140
\qquad\text{y}\qquad
3x+2y=240
\]
se obtiene restando:
\[
2x=100 \Longrightarrow x=50.
\]
Entonces
\[
50+2y=140 \Longrightarrow y=45.
\]
Finalmente, para \(y=10\) en la segunda frontera,
\[
3x+20=240 \Longrightarrow x=\frac{220}{3}.
\]
Los cuatro puntos cumplen todas las restricciones, por lo que los vertices son
\[
(20,10),\qquad (20,60),\qquad (50,45),\qquad \left(\frac{220}{3},10\right).
\]
\textbf{Paso 3. Evaluar la funcion objetivo en todos los vertices.}
Evaluamos el beneficio:
\[
B(20,10)=610,
\quad
B(20,60)=1860,
\quad
B(50,45)=2025,
\]
\[
B\left(\frac{220}{3},10\right)=18\cdot \frac{220}{3}+25\cdot 10=1570.
\]
\textbf{Paso 4. Interpretar.} Una funcion lineal alcanza sus extremos en los vertices de una
region poligonal acotada. El mayor valor, \(2025\), es por tanto el maximo global: el plan
optimo es producir \(50\) cuadernos y \(45\) archivadores.
\end{fullsolution}

\section{C05. Observacion de una torre desde dos puntos}

\begin{contextproblem}{Topografia desde un sendero horizontal}
Desde un sendero horizontal se observa la parte superior de una torre. En el punto \(A\) el angulo
de elevacion es \(28^\circ\). Tras avanzar \(40\) metros hacia la base, en el punto \(B\), el
angulo pasa a ser \(41^\circ\).

Calcula la altura de la torre sobre el terreno horizontal.
\end{contextproblem}

\begin{databox}{Croquis trigonometrico}
\begin{center}
\begin{tikzpicture}[scale=0.95]
  \draw[thick] (0,0) -- (8,0);
  \draw[thick] (8,0) -- (8,4.8);
  \draw[thick] (0,0) -- (8,4.8);
  \draw[thick] (2.4,0) -- (8,4.8);
  \node[below] at (0,0) {\(A\)};
  \node[below] at (2.4,0) {\(B\)};
  \node[below] at (8,0) {Base};
  \node[right] at (8,2.4) {\(h\)};
  \draw (0.9,0) arc[start angle=0,end angle=28,radius=0.9];
  \node at (1.4,0.45) {\scriptsize \(28^\circ\)};
  \draw (3.1,0) arc[start angle=0,end angle=41,radius=0.7];
  \node at (3.55,0.45) {\scriptsize \(41^\circ\)};
\end{tikzpicture}
\end{center}
\end{databox}

\begin{shortanswer}{C05 modelo}
La altura es aproximadamente \(54.8\) metros.
\end{shortanswer}

\begin{fullsolution}{C05 modelo}
\textbf{Paso 1. Definir las distancias.}
Sea \(x\) la distancia desde \(B\) hasta la base de la torre. Entonces desde \(A\) la distancia
horizontal es \(x+40\). Por trigonometria,
\[
\tan 41^\circ=\frac{h}{x},
\qquad
\tan 28^\circ=\frac{h}{x+40}.
\]
\textbf{Paso 2. Expresar la misma altura de dos formas.}
Igualando \(h\),
\[
x\tan 41^\circ=(x+40)\tan 28^\circ.
\]
Despejamos:
\[
x(\tan 41^\circ-\tan 28^\circ)=40\tan 28^\circ,
\]
\[
x=\frac{40\tan 28^\circ}{\tan 41^\circ-\tan 28^\circ}\approx 63.00.
\]
\textbf{Paso 3. Calcular la altura.}
La altura es
\[
h=x\tan 41^\circ\approx 63.00\cdot 0.8693\approx 54.77.
\]
Por tanto, la torre mide aproximadamente \(\SI{54.8}{m}\).

\textbf{Paso 4. Comprobar.} Desde \(B\),
\(54.77/63.00\approx0.8694\approx\tan41^\circ\). Desde \(A\),
\(54.77/103.00\approx0.5317\approx\tan28^\circ\). Ambos angulos se recuperan con el resultado.
\end{fullsolution}

\section{C06. Ruta vectorial de un dron}

\begin{contextproblem}{Tres desplazamientos y una decision de orientacion}
Un dron realiza tres desplazamientos sucesivos, en metros:
\[
\vec u=(6,2),\qquad \vec v=(-1,5),\qquad \vec w=(3,-4).
\]
Se pide:
\begin{enumerate}
  \item hallar el desplazamiento total;
  \item calcular su modulo;
  \item decidir si la trayectoria final esta mas alineada con el vector
  \(\vec a=(2,1)\) o con el vector \(\vec b=(1,3)\).
\end{enumerate}
\end{contextproblem}

\begin{databox}{Esquema del recorrido}
\begin{center}
\begin{tikzpicture}[scale=0.55]
  \draw[->] (-1,0) -- (10,0) node[right] {\(x\)};
  \draw[->] (0,-2) -- (0,8) node[above] {\(y\)};
  \draw[->,thick] (0,0) -- (6,2) node[midway,below] {\(\vec u\)};
  \draw[->,thick] (6,2) -- (5,7) node[midway,left] {\(\vec v\)};
  \draw[->,thick] (5,7) -- (8,3) node[midway,right] {\(\vec w\)};
  \draw[->,very thick,gray] (0,0) -- (8,3) node[midway,above] {\(\vec r\)};
\end{tikzpicture}
\end{center}
\end{databox}

\begin{shortanswer}{C06 modelo}
El desplazamiento total es \(\vec r=(8,3)\), su modulo es \(\sqrt{73}\) y queda mucho mas
alineado con \(\vec a=(2,1)\) que con \(\vec b=(1,3)\).
\end{shortanswer}

\begin{fullsolution}{C06 modelo}
\textbf{Paso 1. Sumar los desplazamientos.}
Sumamos coordenadas:
\[
\vec r=\vec u+\vec v+\vec w=(6,2)+(-1,5)+(3,-4)=(8,3).
\]
\textbf{Paso 2. Calcular el modulo del desplazamiento resultante.}
Su modulo es
\[
\|\vec r\|=\sqrt{8^2+3^2}=\sqrt{73}.
\]
\textbf{Paso 3. Comparar direcciones, no solo productos escalares.}
Para comparar alineacion usamos el producto escalar:
\[
\vec r\cdot \vec a=8\cdot 2+3\cdot 1=19,
\qquad
\vec r\cdot \vec b=8\cdot 1+3\cdot 3=17.
\]
Pero hay que tener en cuenta tambien los modulos:
\[
\cos(\widehat{r,a})=\frac{19}{\sqrt{73}\sqrt{5}}\approx 0.994,
\qquad
\cos(\widehat{r,b})=\frac{17}{\sqrt{73}\sqrt{10}}\approx 0.629.
\]
Los angulos correspondientes son aproximadamente
\[
\widehat{r,a}\approx6.0^\circ,
\qquad
\widehat{r,b}\approx51.0^\circ.
\]
\textbf{Paso 4. Concluir y comprobar.} Como el angulo con \(\vec a\) es mucho menor, la ruta
final queda claramente mas alineada con \(\vec a\). Ademas, el extremo dibujado coincide con la
suma: \((6,2)+(-1,5)+(3,-4)=(8,3)\).
\end{fullsolution}

\section{C07. Triangulo urbano en coordenadas}

\begin{contextproblem}{Parque triangular y rectas notables}
Tres sensores se colocan en los puntos
\[
A(0,0),\qquad B(8,0),\qquad C(2,6).
\]
Se pide:
\begin{enumerate}
  \item calcular el area del triangulo;
  \item hallar la altura trazada desde \(C\) al lado \(AB\);
  \item hallar el ortocentro del triangulo.
\end{enumerate}
\end{contextproblem}

\begin{databox}{Plano de referencia}
\begin{center}
\begin{tikzpicture}[scale=0.62]
  \draw[->] (-1,0) -- (10,0) node[right] {\(x\)};
  \draw[->] (0,-1) -- (0,8) node[above] {\(y\)};
  \draw[thick] (0,0) -- (8,0) -- (2,6) -- cycle;
  \draw[dashed] (2,6) -- (2,0);
  \fill (0,0) circle (2pt) node[below left] {\(A\)};
  \fill (8,0) circle (2pt) node[below right] {\(B\)};
  \fill (2,6) circle (2pt) node[above] {\(C\)};
\end{tikzpicture}
\end{center}
\end{databox}

\begin{shortanswer}{C07 modelo}
El area es \(24\) unidades cuadradas, la altura desde \(C\) es la recta \(x=2\) y el
ortocentro es \(H(2,2)\).
\end{shortanswer}

\begin{fullsolution}{C07 modelo}
\textbf{Paso 1. Calcular el area.}
El lado \(AB\) es horizontal y mide \(8\). La distancia vertical desde \(C\) hasta \(AB\) es
\(6\), asi que
\[
Area=\frac{8\cdot 6}{2}=24.
\]
\textbf{Paso 2. Hallar la altura desde \(C\).}
Como \(AB\) tiene ecuacion \(y=0\), la altura desde \(C(2,6)\) es perpendicular a \(AB\) y por
tanto vertical:
\[
x=2.
\]

\textbf{Paso 3. Hallar una segunda altura.}
Para hallar el ortocentro necesitamos otra altura. La recta \(AC\) tiene pendiente
\[
m_{AC}=\frac{6-0}{2-0}=3.
\]
La altura desde \(B\) debe tener pendiente \(-\frac13\) y pasar por \(B(8,0)\):
\[
y=-\frac13(x-8)=-\frac{x}{3}+\frac{8}{3}.
\]
Intersecando con \(x=2\):
\[
y=-\frac{2}{3}+\frac{8}{3}=2.
\]
Luego el ortocentro es
\[
H(2,2).
\]
\textbf{Paso 4. Comprobar.} El punto \(H\) pertenece a la primera altura porque su abscisa es
\(2\), y pertenece a la segunda porque
\(-2/3+8/3=2\). Por tanto, es la interseccion de ambas alturas.
\end{fullsolution}

\section{C08. Modelo de enfriamiento sencillo}

\begin{contextproblem}{Temperatura estabilizada de un deposito}
La temperatura aproximada de un deposito, en grados, viene dada por
\[
T(t)=2+\frac{12}{t+1},
\]
para \(t\) medido en horas desde que comienza el proceso.

Se pide:
\begin{enumerate}
  \item dar el dominio fisicamente razonable;
  \item estudiar la temperatura a largo plazo;
  \item determinar para que tiempos la temperatura es inferior a \(5^\circ\).
\end{enumerate}
\end{contextproblem}

\begin{databox}{Grafica de apoyo}
\begin{center}
\begin{tikzpicture}
\begin{axis}[
  width=0.78\textwidth,
  height=6cm,
  axis lines=middle,
  xmin=-0.2,
  xmax=10,
  ymin=0,
  ymax=15,
  samples=120,
  grid=both,
  grid style={line width=0.1pt, draw=gray!25},
  major grid style={line width=0.2pt, draw=gray!45},
  xlabel={$t$},
  ylabel={$T(t)$}
]
\addplot[very thick, black, domain=0:10] {2 + 12/(x + 1)};
\addplot[dashed] coordinates {(0,2) (10,2)};
\end{axis}
\end{tikzpicture}
\end{center}
\end{databox}

\begin{shortanswer}{C08 modelo}
En el contexto fisico se toma \(t\geq 0\). La temperatura tiende a \(2^\circ\) a largo plazo y
queda por debajo de \(5^\circ\) para \(t>3\) horas.
\end{shortanswer}

\begin{fullsolution}{C08 modelo}
\textbf{Paso 1. Fijar el dominio del contexto.}
La funcion algebraicamente no esta definida en \(t=-1\), pero en el contexto del problema solo
tiene sentido considerar tiempos desde el inicio, de modo que el dominio fisico es
\[
[0,\infty).
\]
\textbf{Paso 2. Estudiar el comportamiento a largo plazo.}
El limite en el infinito es
\[
\lim_{t\to\infty}\left(2+\frac{12}{t+1}\right)=2.
\]
Eso significa que la temperatura se estabiliza cerca de \(2^\circ\).
Ademas,
\[
T'(t)=-\frac{12}{(t+1)^2}<0\qquad (t\geq0),
\]
por lo que la temperatura decrece durante todo el proceso considerado.

\textbf{Paso 3. Resolver la desigualdad.}
Para encontrar cuando \(T(t)<5\),
\[
2+\frac{12}{t+1}<5
\Longrightarrow
\frac{12}{t+1}<3.
\]
Como \(t\geq 0\), se tiene \(t+1>0\) y se puede multiplicar:
\[
12<3(t+1)
\Longrightarrow
12<3t+3
\Longrightarrow
9<3t
\Longrightarrow
t>3.
\]
\textbf{Paso 4. Comprobar el extremo.} En \(t=3\), \(T(3)=2+12/4=5\), de modo que ese instante
no se incluye. Para cualquier \(t>3\), la temperatura es estrictamente inferior a \(5^\circ\).
\end{fullsolution}

\section{C09. Optimizar el area de un cartel}

\begin{contextproblem}{Panel rectangular apoyado en un muro}
Se quiere construir un panel rectangular apoyado en un muro, de modo que solo hay que vallar
tres lados con una longitud total de \(20\) metros. Si \(x\) es la profundidad y \(y\) el lado
paralelo al muro, calcula las dimensiones que maximizan el area.
\end{contextproblem}

\begin{databox}{Esquema de la optimizacion}
\begin{center}
\begin{tikzpicture}[scale=0.75]
  \draw[very thick] (0,0) -- (9,0);
  \draw[thick] (1,0) rectangle (8,4);
  \node[left] at (1,2) {\(x\)};
  \node[right] at (8,2) {\(x\)};
  \node[below] at (4.5,0) {muro};
  \node[above] at (4.5,4) {\(y\)};
\end{tikzpicture}
\end{center}
\end{databox}

\begin{shortanswer}{C09 modelo}
La funcion area es \(A(x)=20x-2x^2\). El maximo se alcanza en \(x=5\) y entonces \(y=10\).
El area maxima es \(50\) metros cuadrados.
\end{shortanswer}

\begin{fullsolution}{C09 modelo}
\textbf{Paso 1. Traducir la restriccion y fijar el dominio.}
Si se vallan tres lados, la restriccion es
\[
2x+y=20 \Longrightarrow y=20-2x.
\]
Para que ambos lados sean positivos debe cumplirse \(0<x<10\).

\textbf{Paso 2. Construir la funcion objetivo.}
El area del rectangulo es
\[
A(x)=xy=x(20-2x)=20x-2x^2.
\]
\textbf{Paso 3. Localizar y clasificar el punto critico.}
Derivando,
\[
A'(x)=20-4x.
\]
El punto critico cumple
\[
20-4x=0 \Longrightarrow x=5.
\]
La segunda derivada es
\[
A''(x)=-4<0,
\]
luego el punto critico corresponde a un maximo. Como la funcion es una parabola concava y
\(x=5\) pertenece a \((0,10)\), se trata del maximo global en el dominio fisico. Sustituyendo,
\[
y=20-2\cdot 5=10.
\]
El area maxima vale
\[
A(5)=5\cdot 10=50.
\]
\textbf{Paso 4. Interpretar y comprobar.} Las dimensiones optimas son \(5\) metros de
profundidad y \(10\) metros de ancho. Usan exactamente \(2\cdot5+10=20\) metros de valla y
producen un area maxima de \(50\) metros cuadrados.
\end{fullsolution}

\section{C10. Alerta de calidad y decision bajo incertidumbre}

\begin{contextproblem}{Falsos positivos en una linea de produccion}
El \(4\%\) de las piezas de una linea presenta un defecto. Un sistema automatico alerta en el
\(95\%\) de las piezas defectuosas y genera una alerta erronea en el \(3\%\) de las piezas
correctas. Se pide:
\begin{enumerate}
  \item calcular la probabilidad total de alerta;
  \item hallar la probabilidad de que una pieza este defectuosa sabiendo que ha generado alerta;
  \item estimar cuantas revisiones innecesarias se esperan al procesar \(10\,000\) piezas;
  \item explicar por que una alerta no equivale a certeza de defecto.
\end{enumerate}
\end{contextproblem}

\begin{databox}{Arbol de probabilidades}
\begin{center}
\begin{tikzpicture}[x=1cm,y=1cm,>=stealth,every node/.style={font=\small}]
  \node (o) at (0,0) {Pieza};
  \node (d) at (3,1.25) {Defectuosa \(D\)};
  \node (c) at (3,-1.25) {Correcta \(\overline D\)};
  \node (da) at (7,1.8) {Alerta \(A\)};
  \node (dna) at (7,0.7) {Sin alerta};
  \node (ca) at (7,-0.7) {Alerta \(A\)};
  \node (cna) at (7,-1.8) {Sin alerta};
  \draw[->] (o) -- node[above left] {\(0.04\)} (d);
  \draw[->] (o) -- node[below left] {\(0.96\)} (c);
  \draw[->] (d) -- node[above] {\(0.95\)} (da);
  \draw[->] (d) -- node[below] {\(0.05\)} (dna);
  \draw[->] (c) -- node[above] {\(0.03\)} (ca);
  \draw[->] (c) -- node[below] {\(0.97\)} (cna);
\end{tikzpicture}
\end{center}
\end{databox}

\begin{shortanswer}{C10 modelo}
\(P(A)=0.0668\), \(P(D\mid A)\approx0.5689\) y se esperan \(288\) revisiones innecesarias.
\end{shortanswer}

\begin{fullsolution}{C10 modelo}
\textbf{Paso 1. Calcular los dos caminos que producen una alerta.}
Una pieza defectuosa y detectada tiene probabilidad
\[
P(D\cap A)=0.04\cdot0.95=0.038.
\]
Una pieza correcta que produce un falso positivo tiene probabilidad
\[
P(\overline D\cap A)=0.96\cdot0.03=0.0288.
\]

\textbf{Paso 2. Sumar los caminos para obtener la probabilidad total de alerta.}
Por probabilidad total,
\[
P(A)=P(D)P(A\mid D)+P(\overline D)P(A\mid\overline D)
=0.04\cdot0.95+0.96\cdot0.03=0.0668.
\]
\textbf{Paso 3. Invertir la condicion mediante Bayes.}
Aplicando Bayes,
\[
P(D\mid A)=\frac{P(D)P(A\mid D)}{P(A)}
=\frac{0.04\cdot0.95}{0.0668}\approx0.5689.
\]
\textbf{Paso 4. Calcular las revisiones innecesarias.}
Las alertas erroneas esperadas entre \(10\,000\) piezas son
\[
10\,000\cdot0.96\cdot0.03=288.
\]
\textbf{Paso 5. Interpretar y comprobar con frecuencias esperadas.} En \(10\,000\) piezas se
esperan \(380\) alertas verdaderas y \(288\) falsas, es decir, \(668\) alertas en total. La
fraccion \(380/668\approx0.5689\) confirma el resultado de Bayes. La alerta eleva la probabilidad
de defecto del \(4\%\) al \(56.9\%\), pero no equivale a certeza.
\end{fullsolution}

\chapter{Secuencia de Produccion}

\begin{planningbox}{Ruta recomendada de ejecucion}
\begin{enumerate}
  \item Generar una ficha por seccion con contexto, visual, tecnicas y errores previsibles.
  \item Redactar primero los capitulos `C05-C09`, porque son los que mas se benefician de un
  suplemento contextualizado de alta dificultad.
  \item Consolidar despues `C01-C04` con problemas de lectura de datos, algebra y restricciones.
  \item Mantener una plantilla fija: enunciado, recurso visual, preguntas escalonadas,
  respuesta breve, solucion completa y comentario didactico.
  \item Cerrar con una revision visual y una lectura editorial en papel o escala de grises.
\end{enumerate}
\end{planningbox}

\begin{databox}{Recursos que deben prepararse antes de la redaccion masiva}
\begin{itemize}
  \item biblioteca de figuras TikZ para triangulos, vectores, planos, rectas numericas y
  regiones factibles;
  \item tablas reutilizables para datos de produccion, costes, mediciones y lecturas
  experimentales;
  \item convencion estable de dificultad: alta, alta+ y muy alta;
  \item parrilla de control para comprobar que no desaparece la tecnica central del apartado.
\end{itemize}
\end{databox}

\begin{planningbox}{Cierre esperado del suplemento definitivo}
Cuando este proyecto se complete como cuaderno final, el PDF deberia quedar organizado en diez
capitulos, un reto integrador por tema, tres problemas globales de cierre y un apendice de
indices por tecnica, recurso visual y nivel de dificultad.
\end{planningbox}
